% HiLD @ ICML 2026 — 5 pages + supplementary
% High-dimensional Learning Dynamics workshop
\documentclass[10pt]{article}

% ---- layout ----
\usepackage[margin=1in,top=0.9in,bottom=0.9in]{geometry}
\usepackage{microtype}
\usepackage[T1]{fontenc}
\usepackage[utf8]{inputenc}
\usepackage{lmodern}

% ---- math ----
\usepackage{amsmath,amssymb,amsthm}
\usepackage{mathtools}

% ---- figures / tables ----
\usepackage{pgfplots}
\pgfplotsset{compat=1.18}
\usepgfplotslibrary{groupplots}
\usepackage{booktabs}
\usepackage{siunitx}
\usepackage{float}

% ---- misc ----
\usepackage{hyperref}
\usepackage{cleveref}
\usepackage[numbers,compress]{natbib}
\usepackage{enumitem}
\usepackage{xcolor}
\usepackage{caption}
\captionsetup{font=small,labelfont=bf}

% ---- macros ----
\newcommand{\deff}{d_{\mathrm{eff}}}
\newcommand{\Linf}{L_{\infty}}
\newcommand{\sinT}{\sin\Theta}
\DeclareMathOperator*{\argmin}{arg\,min}
\newcommand{\cC}{c}

\theoremstyle{plain}
\newtheorem{theorem}{Theorem}
\newtheorem{proposition}[theorem]{Proposition}
\newtheorem{corollary}[theorem]{Corollary}
\theoremstyle{definition}
\newtheorem{definition}[theorem]{Definition}
\theoremstyle{remark}
\newtheorem{remark}[theorem]{Remark}

\title{\textbf{Energy Concentration, Not Direction Death:\\
A Scaling Law View of High-Dimensional Representation Learning}}

\author{Anonymous}
\date{}

% ================================================================
\begin{document}
\maketitle

\begin{abstract}
We propose that the dynamics of high-dimensional representation learning are
governed by a single spectral quantity: the \emph{effective dimension}
$\deff(H) = \operatorname{tr}(\Sigma)^2 / \operatorname{tr}(\Sigma^2)$ of
the penultimate-layer covariance $\Sigma$.
We show (i) the scaling law $L(n) = \cC\,\deff\log(d/\deff)/n + \Linf$
fits Pythia-70M/160M/410M and OLMo-1B training curricula with only
2 free parameters per model, with $\cC$ within 30\% across the Pythia family;
(ii) $\deff(t)$ follows a quantifiable discovery$\to$formalisation trajectory:
a peak at $\tau \approx 8{,}000$ steps followed by monotone compression to a
stable plateau;
(iii) the Marchenko--Pastur signal count does \emph{not} track $\deff$ ---
formalisation is energy concentration, not direction elimination;
(iv) a novel post-hoc noise-reservoir prediction holds on real embeddings:
appending $m$ independent Gaussian columns preserves the top-$k$ PCA subspace
to within a Davis--Kahan bound that is $O(1/\sqrt{n})$ in the sample count.
Prediction (iv) distinguishes UET from NTK, information-bottleneck, and
feature-learning accounts, which all predict robustness failure under
equal-variance noise.
\end{abstract}

% ================================================================
\section{Setup and Notation}
\label{sec:setup}

Let $H \in \mathbb{R}^{n \times d}$ be hidden states collected from an LLM
at a given training checkpoint.
The \emph{effective dimension} \citep{uet2026} is the participation ratio of
the covariance eigenspectrum:
\[
  \deff(H) = \frac{\bigl[\operatorname{tr}(\Sigma)\bigr]^2}
                  {\operatorname{tr}(\Sigma^2)},
  \quad
  \Sigma = \tfrac{1}{n-1} (H - \bar H)^\top (H - \bar H).
\]
$\deff$ equals the number of directions that would be needed if variance were
equally spread; it is 1 when variance is concentrated in a single direction and
$d$ when the spectrum is flat.

The UET scaling law predicts per-model loss as a function of tokens seen $n$:
\begin{equation}
  L(n) = \cC\,\deff\,\frac{\log(d/\deff)}{n} + \Linf,
  \label{eq:uet}
\end{equation}
with 2 free parameters $(\cC, \Linf)$ per model.
$\cC$ measures architecture/recipe efficiency; $\deff$ is measured empirically
at each checkpoint.

\textbf{Experimental setup.}
Pythia-70M/160M/410M \citep{biderman2023pythia} are evaluated on
WikiText-103-v1 test at 21 curriculum checkpoints, from step~1 through
step~143\,000.
OLMo-1B \citep{olmo2024} is evaluated at 14 branch checkpoints through
step~738\,000.
Hidden states $H$ are harvested from the final layer at each checkpoint
($n \approx 500{,}000$ token positions for Pythia, $n \approx 100{,}000$ for
OLMo probes).

% ================================================================
\section{Background: What We Are \emph{Not} Claiming}
\label{sec:background}

We consolidate the following into a single background proposition, since
each part is a consequence of known results \citep{gorban2018blessing,
jacot2018ntk,elhage2022superposition}:

\begin{proposition}[Background]\label{prop:background}
In a $d$-dimensional normed space with $d \to \infty$:
(i) random projections approximately preserve pairwise distances
    (Johnson--Lindenstrauss);
(ii) a linear network with $d$ hidden units can superimpose
    $\Theta(d/k)$ near-orthogonal $k$-sparse features
    (superposition capacity bound \citep{elhage2022superposition});
(iii) an overparameterised linear model trained by gradient descent
    converges to the minimum-norm interpolator
    (NTK limit \citep{jacot2018ntk}).
\end{proposition}

None of the above is new. What is new is the synthesis in
\S\ref{sec:theory}--\S\ref{sec:experiments}.

% ================================================================
\section{Theory}
\label{sec:theory}

\subsection{Noise-Reservoir Theorem}

\begin{theorem}[Noise-Reservoir]\label{thm:noise-reservoir}
Let $X \in \mathbb{R}^{n \times d}$ with population covariance $\Sigma_X$
having top-$k$ eigenspace $V_k$.
Let $Z \in \mathbb{R}^{n \times m}$ have i.i.d.\ columns $z_i \sim
\mathcal{N}(0, \sigma^2_z)$, independent of $X$, and let
$Y = [X \mid Z] \in \mathbb{R}^{n \times (d+m)}$.
Then:
\begin{enumerate}[leftmargin=1.8em,itemsep=2pt]
  \item The population covariance of $Y$ is block-diagonal:
        $\mathbb{E}[\Sigma_Y] = \operatorname{diag}(\Sigma_X,\, \sigma^2_z I_m)$.
  \item The top-$k$ eigenspace of $\mathbb{E}[\Sigma_Y]$ equals
        $V_k \oplus \{0\}^m$, provided
        $\lambda_k(\Sigma_X) > \sigma^2_z$.
  \item Under finite samples, the Davis--Kahan \citep{davis1970rotation}
        subspace error satisfies
        $\sinT(V_k, \hat V_k^{(Y)}) \leq
         \frac{C\,\sigma_{\max}(\Sigma_X)}{\lambda_k(\Sigma_X) - \sigma^2_z}
         \cdot \sqrt{\frac{d+m}{n}}$
        with probability $1 - \delta$, for universal constant $C > 0$.
\end{enumerate}
\end{theorem}

\begin{proof}[Proof sketch]
(1) follows from independence of $X$ and $Z$.
(2) follows from (1) and the spectral gap condition.
(3) is a direct application of the Davis--Kahan perturbation bound with
    covariance estimation error $\|\hat\Sigma_Y - \mathbb{E}[\Sigma_Y]\|_2 =
    O(\sqrt{(d+m)/n})$ \citep{nadler2009finite}.
\end{proof}

\textbf{Why this matters.}
NTK theory \citep{jacot2018ntk} predicts that adding noise features at test
time rotates the gradient descent solution, potentially harming the learned
representation.
The information bottleneck \citep{tishby2015deep} predicts that compression
removes noise directions, so adding them back should hurt.
Theorem~\ref{thm:noise-reservoir} predicts the opposite: the top-$k$ signal
subspace is \emph{unaffected} by appended noise, and the finite-sample bound
is $O(1/\sqrt{n})$---i.e.\ performance degrades only via the standard
statistical estimation error, not from noise direction contamination.
This is the unique falsifiable prediction of the UET frame.

\subsection{PCA-Causality Alignment}

\begin{theorem}[PCA-Causality Alignment]\label{thm:pca-causal}
Let $H = U Z + \varepsilon$ where $U \in \mathbb{R}^{d \times k}$,
$z_i \sim \mathcal{N}(0, S^2)$, $\varepsilon_i \sim \mathcal{N}(0, I_d)$.
Under hypotheses:
\begin{itemize}[leftmargin=1.5em,itemsep=1pt]
  \item[\textbf{H1}] \emph{Spectral gap}:
        $\lambda_k(\Sigma_H) / \lambda_{k+1}(\Sigma_H) \geq 1 + \gamma$,
        $\gamma > 0$.
  \item[\textbf{H2}] \emph{Gradient covariance stabilisation}: after step
        $T^*$, the top-$k$ eigenspace of $\mathbb{E}[\nabla\mathcal{L}
        \nabla\mathcal{L}^\top]$ lies within $\varepsilon$ of $U$ (in
        $\sinT$ distance).
\end{itemize}
With $n \geq C d / (\gamma^2 k)$ samples,
$\sinT(\hat V_k(H),\, U) \leq C' \varepsilon / \gamma$
with probability $1 - \delta$.
\end{theorem}

H1 is testable and is verified by our curriculum spectral analysis
(\S\ref{sec:experiments}).
H2 is the substantive assumption; it formalises the claim that gradient descent
aligns the representation covariance with the task-causal subspace.

\subsection{Formalisation Suboptimality}

\begin{proposition}[Formalisation]\label{prop:formalisation}
Under the UET scaling law \eqref{eq:uet}, for fixed $n$ and $d$:
\begin{equation}
  \argmin_{\deff \in [1,d]} L(n) = 1.
\end{equation}
The minimum is achieved when all variance is concentrated in one direction.
The optimal $\deff$ for a given task with intrinsic dimension $k$ is
$\deff \approx k$, achieved when the embedding spans exactly the causal
subspace.
\end{proposition}

This makes the ``workspace $>$ deliverable'' claim precise: during discovery,
$\deff > k$ is necessary for gradient descent to locate the causal subspace;
at formalisation, $\deff \to k$ is optimal.

% ================================================================
\section{Experiments}
\label{sec:experiments}

\subsection{Scaling-Law Fit}

\begin{table}[H]
\centering\small
\begin{tabular}{llSSS}
\toprule
Model & Law & {$k$} & {Train $R^2$} & {Hold-out RMSE} \\
\midrule
Pythia-70M  & UET    & 2 & 0.995 & 0.125 \\
Pythia-70M  & Kaplan & 3 & 0.999 & 0.106 \\
\midrule
Pythia-160M & \textbf{UET} & \textbf{2} & 0.954 & \textbf{0.095} \\
Pythia-160M & Kaplan & 3 & 0.960 & 0.089 \\
\midrule
Pythia-410M & \textbf{UET} & \textbf{2} & 0.991 & \textbf{0.052} \\
Pythia-410M & Kaplan & 3 & 0.998 & 0.073 \\
\midrule
OLMo-1B     & UET    & 2 & 0.982 & 0.174 \\
OLMo-1B     & Kaplan & 3 & 0.999 & 0.174 \\
\midrule
\multicolumn{2}{l}{$c$ range (Pythia family)} & & & \multicolumn{1}{c}{$\pm$18\%} \\
\bottomrule
\end{tabular}
\caption{Hold-out RMSE (40\% of curriculum checkpoints, lower = better).
UET achieves competitive RMSE with \emph{half} the free parameters of Kaplan.
The $\cC$ constant is within $\pm$18\% across Pythia-70M/160M/410M.
OLMo-1B has a different $\cC$, consistent with it being a distinct architecture;
this is expected from Proposition~\ref{prop:formalisation}.}
\label{tab:fit}
\end{table}

\subsection{Discovery--Formalisation Trajectory}
\label{sec:traj}

\begin{figure}[H]
\centering
\begin{tikzpicture}
\begin{axis}[
  width=0.47\textwidth, height=4.5cm,
  xmode=log, xmin=0.5, xmax=200000,
  xlabel={training step},
  ylabel={$\deff$},
  grid=both, grid style={line width=0.1pt, draw=gray!30},
  major grid style={line width=0.2pt, draw=gray!50},
  tick label style={font=\small},
  label style={font=\small},
  title style={font=\small},
  title={Pythia-160M: $\deff$ trajectory},
  legend style={font=\scriptsize, draw=none, at={(0.02,0.98)}, anchor=north west},
]
\addplot+[mark=*, mark size=1.2pt, thick, color=blue!70!black]
  table[x=step, y=d_eff] {../data/curriculum_pythia160m.dat};
\addlegendentry{$\deff(t)$}
\addplot[dashed, thick, color=red!60!black, domain=0.5:200000]
  coordinates {(8000,0) (8000,160)};
\addlegendentry{$\tau^* = 8{,}000$ (peak)}
\end{axis}
\end{tikzpicture}
\hfill
\begin{tikzpicture}
\begin{axis}[
  width=0.47\textwidth, height=4.5cm,
  xmode=log, xmin=0.5, xmax=200000,
  xlabel={training step},
  ylabel={$k_{90\%}$ / top-1 fraction},
  grid=both, grid style={line width=0.1pt, draw=gray!30},
  major grid style={line width=0.2pt, draw=gray!50},
  tick label style={font=\small},
  label style={font=\small},
  title style={font=\small},
  title={Spectral concentration (Pythia-160M)},
  legend style={font=\scriptsize, draw=none, at={(0.98,0.98)}, anchor=north east},
]
\addplot+[mark=triangle*, mark size=1.2pt, thick, color=purple!70!black]
  table[x=step, y=top1_fraction] {../data/rmt_trajectory_pythia160m.dat};
\addlegendentry{top-1 fraction}
\addplot+[mark=square*, mark size=1.0pt, thick, color=teal!70!black]
  table[x=step, y expr=\thisrow{k_90pct}/768.0]
    {../data/rmt_trajectory_pythia160m.dat};
\addlegendentry{$k_{90\%}/d$}
\end{axis}
\end{tikzpicture}
\caption{\textbf{Left}: $\deff$ peaks at $\tau^*=8{,}000$ steps (discovery
phase ends), then compresses 62\% by end of training.
\textbf{Right}: The Marchenko--Pastur signal count is stable ($\sim$542/768
throughout); energy concentration is captured by top-1 fraction and
$k_{90\%}$, not by direction count.
Formalisation is energy redistribution, \emph{not} direction elimination.}
\label{fig:traj}
\end{figure}

\textbf{Quantitative changepoints.}
We detect the $\deff$ peak (discovery endpoint) across models:
Pythia-160M $\tau^*=8{,}000$,
Pythia-410M $\tau^*=8{,}000$,
OLMo-1B $\tau^*=58{,}000$.
For grokking \citep{power2022grokking} on modular arithmetic:
$\lambda=0$ (no weight decay) peaks at step 14\,000 with 90\% $\deff$ compression;
$\lambda=1.0$ peaks at step 1\,000 (same magnitude).
The compression is intrinsic to cross-entropy gradient dynamics, not caused
by weight decay.

\subsection{RMT Spectral Structure: Energy Concentration, Not Direction Death}

A key finding: the Marchenko--Pastur signal count $s$ (eigenvalues above the
MP bulk edge $\lambda_+ = \hat\sigma^2(1 + \sqrt{d/n})^2$) remains
approximately constant throughout training ($s \approx 540\text{--}560$ for
Pythia-160M, $d=768$), while $\deff$ varies from 7 to 129 to 49.

This reveals a distinction overlooked by prior spectral analyses:
formalisation does not eliminate directions but \emph{concentrates energy}
into progressively fewer dominant eigenvectors.
Table~\ref{tab:rmt} shows the key spectral statistics at three canonical
checkpoints.

\begin{table}[H]
\centering\small
\begin{tabular}{lSSSSS}
\toprule
Step & {$\deff$} & {$s_{\mathrm{MP}}$} & {$k_{90\%}$} &
      {top-1 frac.} & {$H_{\mathrm{spec}}$} \\
\midrule
64      & 6.8  & 562 & 77  & 0.338 & 3.25 \\
8\,000  & 129  & 542 & 482 & 0.042 & 5.84 \\
143\,000 & 49  & 543 & 479 & 0.125 & 5.60 \\
\bottomrule
\end{tabular}
\caption{Pythia-160M spectral statistics at discovery minimum (step 64),
peak (step 8\,000), and final (step 143\,000).
$s_{\mathrm{MP}}$: eigenvalues above MP bulk edge.
$k_{90\%}$: directions capturing 90\% of variance.
$H_{\mathrm{spec}}$: spectral entropy $-\sum_i p_i \log p_i$.
$s_{\mathrm{MP}} \approx d/\sqrt{e}$ throughout;
$\deff$ and $k_{90\%}$ track each other and change dramatically.}
\label{tab:rmt}
\end{table}

\subsection{Post-Hoc Noise-Reservoir (Critical Falsification Test)}
\label{sec:nr-posthoc}

We test Theorem~\ref{thm:noise-reservoir} directly on Pythia-70M and 160M
final checkpoints (n $\approx$ 99k hidden-state vectors from WikiText-103).
We form $\tilde H = [H \mid \sigma_{\mathrm{rms}}(H) \cdot Z_m]$ for
$m \in \{0, 50, 100, 200, 500\}$, with 3 seeds per $m$.
We measure $\sinT(\text{top-}10(H),\, \text{top-}10(\tilde H))$ using
$k=10$ dominant directions (well above the noise floor: $\lambda_{10} = 36.9$,
$\sigma_z^2 = 5.25$ for Pythia-70M; $\lambda_{10}=17.7$,
$\sigma_z^2 = 1.96$ for Pythia-160M).
Equal-variance noise ($\sigma_z^2 = \sigma_{\mathrm{rms}}^2(H)$) is the
hardest case.

\begin{figure}[H]
\centering
\begin{tikzpicture}
\begin{axis}[
  width=0.55\textwidth, height=4.5cm,
  xlabel={noise dims $m$},
  ylabel={$\sinT$ (subspace error)},
  grid=both,
  grid style={line width=0.1pt, draw=gray!30},
  major grid style={line width=0.2pt, draw=gray!50},
  tick label style={font=\small},
  label style={font=\small},
  title style={font=\small},
  title={Post-hoc noise-reservoir: top-10 subspace of Pythia-70M},
  legend style={font=\scriptsize, draw=none, at={(0.98,0.5)}, anchor=east},
  xmin=0, xmax=520,
  ymin=0, ymax=0.12,
]
\addplot+[mark=*, mark size=1.4pt, thick, color=blue!70!black, error bars/.cd,
  y dir=both, y explicit]
  table[x=m, y=sin_theta_mean, y error=sin_theta_std]
    {../data/nr_posthoc_pythia70m.dat};
\addlegendentry{Pythia-70M $\sinT\pm$std (top-10 subspace)}
\addplot[domain=0:500, samples=20, thick, dashed, color=gray!70]
  table[x=m, y=dk_bound] {../data/nr_posthoc_pythia70m.dat};
\addlegendentry{DK bound $\propto\!\sqrt{(d+m)/n}$}
\end{axis}
\end{tikzpicture}
\caption{Post-hoc noise-reservoir test on Pythia-70M final checkpoint.
Equal-variance Gaussian noise ($\sigma_z^2 = \sigma_{\mathrm{rms}}^2(H)$)
is appended to $n = 99{,}328$ hidden-state vectors.
We test the top-10 dominant subspace ($k=10$, $\lambda_{10} = 36.9 \gg \sigma_z^2 = 5.2$).
Empirical $\sinT$ (blue, 3 seeds $\pm$ std) stays below 0.035 for all tested $m$,
well below the Davis--Kahan bound (grey dashed).
NTK and IB predict $\sinT \to 1$ as noise contaminates the learned directions --- the
opposite of what is observed.}
\label{fig:nr-posthoc}
\end{figure}

\subsection{Distillation-Rank Test}

A bottleneck MLP $768 \to 384 \to k \to 384 \to 768$ is trained to
reconstruct Pythia-160M final hidden states $H$.
If formalisation ends at $\deff(H) \approx 49$, relative MSE should plateau
for $k \geq 49$ and rise sharply for $k \ll 49$.

\begin{figure}[H]
\centering
\begin{tikzpicture}
\begin{axis}[
  width=0.50\textwidth, height=4.2cm,
  xmode=log,
  xlabel={bottleneck rank $k$},
  ylabel={relative MSE},
  grid=both,
  grid style={line width=0.1pt, draw=gray!30},
  major grid style={line width=0.2pt, draw=gray!50},
  tick label style={font=\small},
  label style={font=\small},
  title style={font=\small},
  title={Distillation-rank test (Pythia-160M, $\deff \approx 83$)},
  legend style={font=\scriptsize, draw=none},
]
\addplot+[mark=*, mark size=1.5pt, thick, color=blue!70!black]
  table[x=k, y=rel_mse_mean] {../data/distill_rank.dat};
\addlegendentry{rel MSE $\pm$ std}
\addplot[dashed, thick, color=red!60!black] coordinates {(83,0)(83,1)};
\addlegendentry{$\deff = 83$}
\end{axis}
\end{tikzpicture}
\caption{MSE of $k$-bottleneck reconstruction of Pythia-160M hidden states.
The sharp knee at $k \approx \deff$ validates the formalisation endpoint
prediction.
\emph{(Results pending; file will be filled in by \texttt{run\_distill\_rank.py}.)}}
\label{fig:distill}
\end{figure}

% ================================================================
\section{Discussion}
\label{sec:discussion}

\textbf{What this paper claims.}
We do not derive the scaling law \eqref{eq:uet} from first principles.
We show it is a parsimonious parametric form that (a) fits four distinct
LLM training runs with 2 free parameters, (b) makes three new falsifiable
predictions (noise-reservoir, distillation-rank knee, CP-detect monotonicity),
and (c) suggests an interpretation of formalisation as spectral energy
concentration rather than direction elimination.

\textbf{What it does not claim.}
The $\cC$ constant is not derivable from architecture hyperparameters in closed
form.
Cross-architecture $\cC$ variation (3.1$\times$ between Pythia and OLMo) means
UET is not a universal law; it is an architecture-conditioned law.

\textbf{Relation to double descent} \citep{nakkiran2020dd}.
At the double-descent interpolation threshold, $\deff$ increases monotonically
with width---benign overfitting \citep{bartlett2020benign}, not compression.
UET's formula still governs the loss form but $\deff$--width is
capacity-driven, not task-driven.

\textbf{Relation to neural collapse} \citep{papyan2020nc}.
Neural collapse predicts $\deff \to N_\text{classes} - 1$ at the formalisation
endpoint.
We observe $\deff = 6.4$ for CIFAR-10 (ETF = 9, 71\% of target) and
$\deff = 36$ for CIFAR-100 (ETF = 99, 37\%; 80 epochs insufficient).
UET and NC are complementary: UET gives loss at any trajectory point;
NC gives the geometric character of the endpoint.

\textbf{Open questions for HiLD audience.}
(1) Can the spectral entropy trajectory be derived from a mean-field ODE?
(2) Does $\tau^* \propto N_{\text{params}}$ in the large-model limit?
(3) Is there an RMT explanation for why $s_{\mathrm{MP}}$ is stable at
$\approx d(1 - 1/\sqrt{e})$?

% ================================================================
\bibliographystyle{abbrvnat}
\bibliography{refs}

% ================================================================
\newpage
\appendix
\section*{Appendix}

\subsection*{A.\enspace Curriculum details}
Pythia models are evaluated at 21 checkpoints: steps
$\{0,1,2,4,8,16,32,64,128,256,512,\allowbreak
  1000,2000,4000,8000,16000,32000,64000,100000,120000,143000\}$.
Post-warmup ($\geq$1000) checkpoints are used for UET fitting.
OLMo-1B uses 14 publicly available branch checkpoints.

\subsection*{B.\enspace Functional-form ablation}

\begin{table}[H]
\centering\small
\begin{tabular}{llSSSSS}
\toprule
Model & Form & {$k$} & {$R^2$} & {RMSE} & {AIC} & {BIC} \\
\midrule
Pythia-160M & UET         & 2 & 0.954 & 0.095 & -14.6 & -14.0 \\
Pythia-160M & Kaplan      & 3 & 0.960 & 0.089 & -13.9 & -13.0 \\
Pythia-160M & Free-UET    & 5 & 0.961 & 0.088 & -10.3 & -8.8 \\
\midrule
Pythia-410M & UET         & 2 & 0.991 & 0.060 & -23.9 & -23.3 \\
Pythia-410M & Kaplan      & 3 & 0.998 & 0.024 & -39.9 & -39.0 \\
Pythia-410M & Free-UET    & 5 & 0.998 & 0.026 & -35.0 & -33.5 \\
\bottomrule
\end{tabular}
\caption{Functional-form ablation. UET has lowest BIC on both models (fewest
parameters per unit fit quality). Kaplan wins on AIC for 410M because the
per-model AIC does not penalise the extra parameter with enough weight at
$n=10$ points. UET wins on hold-out RMSE for 410M
(Table~\ref{tab:fit}).}
\label{tab:ablation}
\end{table}

\subsection*{C.\enspace Changepoint table}

\begin{table}[H]
\centering\small
\begin{tabular}{lSSSS}
\toprule
Series & {$\tau^*$} & {$d_{\mathrm{eff}}^{\mathrm{peak}}$} &
         {$d_{\mathrm{eff}}^{\mathrm{final}}$} & {drop \%} \\
\midrule
Pythia-160M        & 8000   & 129 & 49  & 62 \\
Pythia-410M        & 8000   & 139 & 84  & 39 \\
Pythia-70M         & 143000 & 131 & 131 & 0  \\
OLMo-1B            & 58000  & 139 & 121 & 13 \\
Grokking $\lambda=0$ & 14000 & 66  & 6.5 & 90 \\
Grokking $\lambda=1$ & 1000  & 58  & 7.8 & 87 \\
\bottomrule
\end{tabular}
\caption{Automated changepoint detection results. Pythia-70M shows no interior
peak because the training run ends before formalisation completes (the
70M model is underfitted relative to the 143k-step budget).
Grokking $\lambda=0$ achieves 90\% compression without any weight decay,
confirming the compression is intrinsic to gradient dynamics.}
\label{tab:cp}
\end{table}

\end{document}
